% !TEX root = ../main.tex

\section{Endoodwracalny kwantowy silnik Otta}
Kwantowy cykl Otta jest serią naprzemiennych przemian o zerowej pracy i zerowym cieple. Zgodnie z rysunkiem \ref{fig:pv_silnik_otto} zmiany energii po wycałkowaniu mają postać:
\begin{align*}
	Q_{1-2} &= 0, &
	W_{1-2} &= U\pab{T_2,\omega_2} - U\pab{T_1,\omega_1}, \\
	Q_{2-3} &= U\pab{T_3,\omega_2} - U\pab{T_2,\omega_2}, &
	W_{2-3} &= 0, \\
	Q_{3-4} &= 0, &
	W_{3-4} &= U\pab{T_4,\omega_1} - U\pab{T_3,\omega_2}, \\
	Q_{4-1} &= U\pab{T_1,\omega_1} - U\pab{T_4,\omega_2}, &
	W_{4-1} &= 0.
\end{align*}
Parametrem kontrolowanym przez nas jest \(\omega\). Rozpatrujemy przypadek w granicy niskiego oddziaływania z otoczeniem. W tym przypadku stan stacjonarny jest stanem równowagi termodynamicznej opisywanym rozkładem Boltzmanna \cite{endo_otto}.

Czynnikiem roboczym w rozpatrywanym cyklu jest oscylator harminiczny opisany w dodatku \ref{sec:osc_harm}. Ważniejsze wielkości termodynamiczne zawarto we wzorach \eqref{eq:osc_harm_Z} -- \eqref{eq:osc_harm_S}.
Funcja \(S\) jest zależna jedynie od \(\frac{\omega}{T}\) i możemy zapisać:
\begin{equation*}
	\tilde{S}\pab{\frac{\omega}{T}} \overset{def}{=} S\pab{T,\omega}.
\end{equation*}
Jest to funkcja monotoniczna, więc jest ona też odwracalna. Oznacza to, że jeśli wartości funkcji są równe to również argumenty tej funcji są równe.

Podczas przemiany izentropowej entropia na początku i końcu przemianu jest równa, stąd mamy:
\begin{equation}\label{eq:endo_otto_omegaT}
\begin{split}
	S\pab{T_1,\omega_1} = S\pab{T_2,\omega_2} \quad\Rightarrow\quad
	\tilde{S}\pab{\frac{\omega_1}{T_1}} = \tilde{S}\pab{\frac{\omega_2}{T_2}} \quad\Rightarrow\quad
	\frac{\omega_1}{T_1} = \frac{\omega_2}{T_2}, \\
	S\pab{T_3,\omega_2} = S\pab{T_4,\omega_1} \quad\Rightarrow\quad
	\tilde{S}\pab{\frac{\omega_2}{T_3}} = \tilde{S}\pab{\frac{\omega_1}{T_4}} \quad\Rightarrow\quad
	\frac{\omega_2}{T_3} = \frac{\omega_1}{T_4}.
\end{split}
\end{equation}
Prace i ciepła w odpowiednich przemianach są równe:
\begin{gather*}
	W_{1-2} = \frac{\hbar}{2}\pab{\omega_2 - \omega_1} \ctgh\pab{\frac{\hbar\omega_1}{2k_B T_1}}, \\
	W_{3-4} = \frac{\hbar}{2}\pab{\omega_1 - \omega_2} \ctgh\pab{\frac{\hbar\omega_2}{2k_B T_3}}, \\
	Q_{2-3} = \frac{\hbar\omega_2}{2} \ctgh\pab{\frac{\hbar\omega_2}{2k_B T_3}} - \frac{\hbar\omega_2}{2} \ctgh\pab{\frac{\hbar\omega_1}{2k_B T_1}}.
\end{gather*}
Wprowadzając wielkość pomocniczą \(\kappa = \frac{\omega_1}{\omega_2}\) znajdujemy sprawność silnika:
\begin{equation}\label{eq:sprawnosc_kwantowy_otto}
	\eta = - \frac{W_{1-2} + W_{3-4}}{Q_{2-3}} = 1 - \kappa.
\end{equation}
Jest to wzór analogiczny do wzoru \eqref{eq:parametry_otto_eta} w tym, że jest on niezależny od temperatur zbiorników.

%Przyjmujemy, że ciepło dostarczone i oddane spełnia równanie Newtona, co można zapisać \cite{endo_otto}:
%\begin{align}
%	\odv{T}{t} &= -\alpha_h \pab{T - T_h}, &
%	T\pab{t=0} &= T_2, &
%	T\pab{t=\tau_h} &= T_3, \\
%	\odv{T}{t} &= -\alpha_c \pab{T - T_c}, &
%	T\pab{t=0} &= T_4, &
%	T\pab{t=\tau_c} &= T_1,
%\end{align}
%gdzie \(T_h\) oraz \(T_c\) to odpowiednio temperatury zbiornika ciepłego i zimnego, które nie są osiągane przez układ. Rozwiązując równania różniczkowe z zadanymi warunkami brzegowymi otrzymujemy:
%\begin{align*}
%	T_3 - T_h &= \pab{T_2 - T_h} e^{-\alpha_h \tau_h}, &
%	T_1 - T_c &= \pab{T_4 - T_c} e^{-\alpha_c \tau_c}, \\
%	T_2 - T_h &= \pab{T_3 - T_h} e^{\alpha_h \tau_h}, &
%	T_4 - T_c &= \pab{T_1 - T_c} e^{\alpha_c \tau_c}.
%\end{align*}
%Możemy je przepisać:
%\begin{align*}
%	T_3 - T_2 e^{-\alpha_h \tau_h} &= T_h \pab{1 - e^{-\alpha_h \tau_h}}, &
%	T_1 - T_4 e^{-\alpha_c \tau_c} &= T_c \pab{1 - e^{-\alpha_c \tau_c}}, \\
%	T_3 e^{\alpha_h \tau_h} - T_2 &= T_h \pab{e^{\alpha_h \tau_h} - 1}, &
%	T_1 e^{\alpha_c \tau_c} - T_4 &= T_c \pab{e^{\alpha_c \tau_c} - 1}.
%\end{align*}
%Z powyższych wzorów oraz \eqref{eq:endo_otto_omegaT} możemy wyprowadzić wzory na temperaturę zależne od \(\tau_h\), \(\tau_c\) oraz \(\kappa\):
%\begin{equation}\label{eq:endo_otto_T}
%\begin{split}
%	T_1 &= \frac{
%	T_c e^{\alpha_h \tau_h} \pab{e^{\alpha_c \tau_c} - 1} + \kappa T_h \pab{e^{\alpha_h \tau_h} - 1}
%	}{
%	e^{\alpha_h \tau_h + \alpha_c \tau_c} - 1
%	}, \\
%	T_2 &= \frac{
%	T_c e^{\alpha_h \tau_h} \pab{e^{\alpha_c \tau_c} - 1} + \kappa T_h \pab{e^{\alpha_h \tau_h} - 1}
%	}{
%	\kappa \pab{e^{\alpha_h \tau_h + \alpha_c \tau_c} - 1}
%	}, \\
%	T_3 &= \frac{
%	T_c \pab{e^{\alpha_c \tau_c} - 1} + \kappa T_h e^{\alpha_c \tau_c} \pab{e^{\alpha_h \tau_h} - 1}
%	}{
%	\kappa \pab{e^{\alpha_h \tau_h + \alpha_c \tau_c} - 1}
%	}, \\
%	T_4 &= \frac{
%	T_c \pab{e^{\alpha_c \tau_c} - 1} + \kappa T_h e^{\alpha_c \tau_c} \pab{e^{\alpha_h \tau_h} - 1}
%	}{
%	e^{\alpha_h \tau_h + \alpha_c \tau_c} - 1
%	}.
%\end{split}
%\end{equation}

Moc silnika wyrażona jest wzorem:
\begin{equation*}
	P = \frac{-W}{\zeta\pab{\tau_h + \tau_c}} = \frac{\hbar\omega_2\pab{1-\kappa}}{2\zeta\pab{\tau_h + \tau_c}} \bab{\ctgh\pab{\frac{\hbar\omega_2}{2k_B T_3}} - \ctgh\pab{\frac{\hbar\omega_2}{2k_B T_2}}}.
\end{equation*}
Wykorzystując przekształcenie:
\begin{multline*}
	\ctgh\pab{x} - \ctgh\pab{y} =
	\frac{\cosh\pab{x}}{\sinh\pab{x}} - \frac{\cosh\pab{y}}{\sinh\pab{y}} =
	\frac{\sinh\pab{y}\cosh\pab{x} - \cosh\pab{y}\sinh\pab{x}}{\sinh\pab{x}\sinh\pab{y}} = \\ =
	\frac{\sinh\pab{y-x}}{\sinh\pab{x}\sinh\pab{y}} =
	\csch\pab{y}\csch\pab{x}\sinh\pab{y-x},
\end{multline*}
możemy alternatywnie zapisać ten wzór jako:
\begin{equation*}
	P = \frac{\hbar\omega_2\pab{1-\kappa}}{2\zeta\pab{\tau_h + \tau_c}} \csch\pab{\frac{\hbar\omega_2}{2k_B T_2}} \csch\pab{\frac{\hbar\omega_2}{2k_B T_3}} \sinh\bab{\frac{\hbar\omega_2}{2k_B}\pab{\inv{T_2}-\inv{T_3}}}.
\end{equation*}
Podstawiając \eqref{eq:endo_otto_T} otrzymujemy zależność mocy od \(\tau_h\), \(\tau_c\) oraz \(\kappa\) \cite{endo_otto}:
\begin{multline}
	P\pab{\tau_h,\tau_c,\kappa} =
	\frac{\hbar\omega_2\pab{1-\kappa}}{2\zeta\pab{\tau_h + \tau_c}}
	\csch\pab{\frac{\hbar\omega_2 \kappa}{2k_B}
	\frac{
		e^{\alpha_h \tau_h + \alpha_c \tau_c} - 1
	}{
		T_c e^{\alpha_h \tau_h} \pab{e^{\alpha_c \tau_c} - 1} + \kappa T_h \pab{e^{\alpha_h \tau_h} - 1}
	}} \\
	\csch\pab{\frac{\hbar\omega_2 \kappa}{2k_B}
	\frac{
		e^{\alpha_h \tau_h + \alpha_c \tau_c} - 1
	}{
		T_c \pab{e^{\alpha_c \tau_c} - 1} + \kappa T_h e^{\alpha_c \tau_c} \pab{e^{\alpha_h \tau_h} - 1}
	}} \\
	\sinh\pab{\frac{\hbar\omega_2 \kappa}{2k_B}
	\frac{
	\pab{\kappa T_h - T_c}\pab{e^{\alpha_h \tau_h + \alpha_c \tau_c} - 1}\pab{e^{\alpha_h \tau_h} - 1}\pab{e^{\alpha_c \tau_c} - 1}
	}{
	\bab{T_c \pab{e^{\alpha_c \tau_c} - 1} + \kappa T_h e^{\alpha_c \tau_c} \pab{e^{\alpha_h \tau_h} - 1}}
	\bab{T_c e^{\alpha_h \tau_h} \pab{e^{\alpha_c \tau_c} - 1} + \kappa T_h \pab{e^{\alpha_h \tau_h} - 1}}
	}}.
\end{multline}
Wynik ten jest nieseparowalny i zbyt skompliwowany, by go maksymalizować analitycznie. Należy go rozwiązać numerycznie.

%W klasycznym silniku Stirlinga otrzymujemy analogiczne rezultaty \eqref{eq:otto_stosunek_T}. Temperatury są również równe \eqref{eq:endo_otto_T}. Sprawność i moc tego silnika wyrażona jest wzorem:
%\begin{equation}
%	\eta = 1 - \kappa \label{eq:otto_eta},
%\end{equation}
%\begin{multline}
%	P\pab{\kappa,\tau_h,\tau_c} = \frac{n c_V \pab{T_1 - T_2 + T_3 - T_4}}{\zeta\pab{\tau_h + \tau_c}} = \\ =
%	\frac{2 n c_V}{\zeta} \cdot \frac{1-\kappa}{\kappa}\pab{\kappa T_h - T_c} \cdot \inv{\tau_h + \tau_c}\sinh\pab{\frac{\alpha_h \tau_h}{2}}\sinh\pab{\frac{\alpha_c \tau_c}{2}}\csch\pab{\frac{\alpha_h \tau_h + \alpha_c \tau_c}{2}} \label{eq:otto_P}.
%\end{multline}
%Wyrażenie na moc zależy jedynie od \(\kappa\), natomiast na moc można rozseparować na czynnik zależny od \(\kappa\) i od pozostałych zmiennych. Aby znaleźć sprawność w punkcie maksymalnej mocy wystarczy znaleźć maksimum funkcji zależnej od \(\kappa\):
%\begin{gather*}
%	f\pab{\kappa} = \frac{1-\kappa}{\kappa}\pab{\kappa T_h - T_c}, \\
%	\odv{f}{\kappa} = \inv{\kappa_{max}^2}T_c - T_h = 0, \\
%	\kappa_{max} = \sqrt{\frac{T_c}{T_h}}, \\
%	\eta\pab{\eta_{max}} = 1 - \sqrt{\frac{T_c}{T_h}}.
%\end{gather*}
%Jest to wynik identyczny do sprawności cyklu Carnota w punkcie maksymalnej mocy \eqref{eq:endo_carnot_eta}.
